% TEMPLATE DE SUBMISSAO DE RESUMO PARA APRESENTACAO DE TRABALHO
% NA II MOSTRA DA PÓS-GRADUAÇÃO NA SEMANA DE MATEMÁTICA APLICADA - SEMAP

% BASEADO NOS TEMPLATES OFICIAIS DA SBMAC PARA O CNMAC 
% PARA MAIS INSTRUCOES SOBRE FIGURAS, TABELAS E EQUACOES,
% BAIXE O ARQUIVO DISPONIVEL EM
% http://www.cnmac.org.br/Template-Categoria_3-SBMAC.zip

\documentclass[a4]{report}

\usepackage[brazil]{babel}      % para texto em Português
%\usepackage[english]{babel}    % para texto em Inglês

%\usepackage[latin1]{inputenc}   % para acentuação em Português OU
\usepackage[utf8x]{inputenc}   % para acentuação em Português com o uso do Unicode, 
% mude a codificação desse template para utf-8

%%
%% POR FAVOR, EVITE MUDAR A CONFIGURAÇÃO DESTE TEMPLATE
%% ADICIONE APENAS OS PACOTES ESSENCIAIS PARA SEU TRABALHO
%%
\usepackage{graphics}
\usepackage{subfigure}
\usepackage{graphicx}
\usepackage{epsfig}
\usepackage[centertags]{amsmath}
\usepackage{graphicx,indentfirst,amsmath,amsfonts,amssymb,amsthm,newlfont}
\usepackage{longtable}
\usepackage{cite}
\usepackage[usenames,dvipsnames]{color}
\usepackage{lipsum}


\begin{document}

%********************************************************
% TÍTULO DO TRABALHO
\title{Instruções para Submissão de Trabalhos na II SEMAP}

% PRIMEIRO AUTOR E EMAIL DE CONTATO
\author{
 {\large  Johann C. F. Gauss }\thanks{dekan@math-cs.uni-goettingen.de}\\
 {\small Departamento de Matemática, University of Göttingen, Göttingen, Alemanha} \\
% % SEGUNDO AUTOR E EMAIL DE CONTATO
%{\large Mateus Bernardes}\thanks{mbernardes@uftpr.edu.br} \\
%{\small Departamento Acadêmico de Matemática, UTFPR, Curitiba, PR} \\
%% TERCEIRO AUTOR E EMAIL DE CONTATO
% {\large Igor Leite Freire}\thanks{igor.freire@ufabc.edu.br}\\
% {\small Centro de Matemática, Computação e Cognição, UFABC}  
 }

\criartitulo

% A CRIAÇÃO DE SESSÕES É LIVRE, DESDE QUE RESPEITADO O LIMITE DE 3 PÁGINAS.

\section{ Introdução }

\lipsum[2-3]

\section{Desenvolvimento}

\lipsum[5-8]

\section{Conclusões}
\lipsum[15]

\section*{Agradecimentos}
Agradecemos à Deutsche Forschungsgemeinschaft pelo apoio financeiro.

% POR FAVOR, MANTENHA A BIBLIOGRAFIA NO PADRÃO RECOMENDADO

\begin{thebibliography}{00}

\bibitem{Boldrini} J. L. Boldrini, S. I. R. Costa, V. R. Ribeiro, and H. G. Wetzler. {\it Álgebra Linear 
e Aplicações}. Harper-Row, São Paulo, 1987.

\bibitem{Cuminato}
J. A. Cuminato and V. Ruas, Unification of distance inequalities for linear variational problems, 
{\it Comp. Appl. Math.}, 2014. DOI: 10.1007/s40314-014-0163-6.

\bibitem{Diniz1}
G. L. Diniz, A mudança no habitat de populações de peixes: de rio a represa -- o modelo matemático, Dissertação de Mestrado em Matemática Aplicada, Unicamp, (1994).

\bibitem{Jafelice} R. M. Jafelice, L. C. Barros and R. C. Bassanezi. Study of the dynamics of HIV under 
treatment considering fuzzy delay, {\it Comp. Appl. Math.}, 33:45--61, 2014.

\bibitem{Santos} I. L. D. Santos e G. N. Silva. Uma classe de problemas de controle ótimo 
em escalas temporais, {\it Proceeding Series of the Brazilian Society of Computational and 
Applied Mathematics}, volume 1, 2013. DOI: 10.5540/03.2013.001.01.0177.

\end{thebibliography}
\end{document}